\begin{textsample}{NEG Dim 1 – 2009\_c4 – Score -8.00 – 2009\_c4/enade\_2009\_teatro\_7552cca521.txt}  \label{ex:f1_neg_027}
Capítulo 4 Percepção sobre a Prova As percepções dos estudantes sobre a prova aplicada no ENADE/2009 na área de TEATRO foram mensuradas por meio de nove itens que avaliaram desde o grau de dificuldade da prova até o tempo gasto para concluí -la.
As questões foram analisadas separando -se concluintes e ingressantes, e as percepções sobre a prova foram relacionadas com o desempenho dos estudantes e com a Região Geográfica de \textbf{localização} do curso6.
O desempenho dos estudantes dentro de \textbf{cada} Área/Subárea foi classificado em quartos ( 1/4 ) para os concluintes e para os ingressantes em separado ( ver Figura 1 para o \textbf{esquema} de representação gráfica destas grandezas ).
As notas finais foram colocadas em \textbf{ordem} ascendente e no primeiro quarto foram colocadas as 25 \% notas mais baixas, limitadas superiormente pelo percentil 25 que vem a ser o primeiro quartil.
No segundo quarto, as 25 \% seguintes, limitadas superiormente pelo percentil 50 que vem a ser a mediana.
No terceiro quarto, as 25 \% notas após as encontradas no segundo quarto, ou seja, as primeiras 25 \% acima da mediana, limitadas superiormente pelo percentil 75 que vem a ser o terceiro quartil.
No último quarto foram alocadas as 25 \% maiores notas.
Potencialmente, poderia existir uma relação entre a nota do estudante e a sua percepção sobre a prova.
Calculando as médias das respostas segundo os quartos de notas, seria possível reconhecer a existência ( ou não ) de algum gradiente nas respostas, por exemplo, uma melhor nota correspondendo a uma percepção menos negativa.
Este gradiente seria traduzido, por exemplo, no gráfico 4.1a, pela diminuição das áreas laranja e vermelha ( Difícil e Muito difícil ) com o aumento do desempenho ( indo do primeiro ao último quarto ). 6 Curso, neste relatório, refere -se à unidade de análise para o Conceito ENADE e é caracterizado pela combinação de IES, Área/Subárea ( quando pertinente ) e município de habilitação.
A medida tradicional de “ grau de discriminação ” comparou os resultados dos estudantes do primeiro quarto ( 25 \% piores estudantes ) aos do último quarto ( 25 \% melhores ).
Esta comparação, no gráfico disponibilizado, foi idêntica à comparação da primeira barra ( primeiro quarto ) com a última barra ( último quarto ) de \textbf{cada} grupo ingressante ou concluinte nos gráficos “ segundo quartos de nota final ” ( gráficos com números ímpares neste capítulo ).
Por outro lado, a disponibilização das informações segundo os quartos de nota permitiu verificar o comportamento de dois outros grupos intermediários que na análise somente dos extremos tem que ser suposto monotônico.
A seguir, serão apresentados os principais resultados relativos aos nove itens avaliados. 4.1 Grau de dificuldade da prova 4.1.1 Formação Geral Os Gráficos 4.1a, 4.1b e 4.1c apresentam a distribuição das respostas à questão “ Qual o grau de dificuldade desta prova na parte de Formação Geral? ”, segundo quartos de nota final e grupo de estudantes concluintes e ingressantes que participaram do ENADE/2009 na área de TEATRO.
No total de concluintes, 12,7 \% declararam que a prova na parte da Formação Geral foi Difícil/Muito difícil, contra um valor maior para os ingressantes, 20,5 \%.
Tanto entre os concluintes, como entre os ingressantes, a proporção de indivíduos que classificaram esta parte como Difícil/Muito difícil diminui com o desempenho do estudante.
Por outro lado, a proporção que declarou a parte Fácil/Muito fácil foi ligeiramente maior entre os concluintes ( 27,9 \% ) do que entre os ingressantes ( 22,9 \% ).
Tanto entre os concluintes como entre os ingressantes esta proporção aumenta com o desempenho.
Entre os concluintes, 59,4 \% classificaram esta parte com um nível Médio de dificuldade.
A proporção entre os ingressantes foi ligeiramente menor, 56,6 \%.
Não se observou uma variação monotônica nem entre os concluintes nem entre os ingressantes para esta proporção em relação ao desempenho do estudante.
Os Gráficos 4.2a e 4.2b apresentam a distribuição das respostas à questão “ Qual o grau de dificuldade desta prova na parte de Formação Geral? ”, segundo a Região onde o curso \textbf{está} situado para os estudantes concluintes e ingressantes que participaram do ENADE/2009 na área de TEATRO.
Cumpre \textbf{lembrar} que a Região Norte não apresentou nenhum estudante concluinte.
Entre os concluintes, a maior proporção de estudantes que considerou a parte de Formação Geral da prova como de dificuldade média entre todas as Regiões brasileiras, foram os da Região Sudeste com 61,7 \%, já entre os ingressantes, foram os estudantes da Região Norte com 59,5 \%.
A Região com as menores proporções desta resposta foram respectivamente para os concluintes e ingressantes, a Centro-Oeste ( 55,2 \% e 51,5 \% ). 4.1.2 Componente Específico Os Gráficos 4.3a, 4.3b e 4.3c apresentam a distribuição das respostas à questão “ Qual o grau de dificuldade desta prova na parte do Componente Específico? ”, segundo quartos de nota final e grupo de estudantes concluintes e ingressantes que participaram do ENADE/2009 na área de TEATRO.
Em geral, para as desagregações consideradas, os estudantes classificaram o Componente Específico como mais difícil que a Formação Geral.
No total de concluintes, 26,9 \% declararam que a prova na parte do Componente Específico foi Difícil/Muito difícil, sendo um valor muito maior para os ingressantes, 48,9 \%.
A proporção que declarou Difícil/Muito difícil entre os concluintes é declinante como função do desempenho : indo de 37,8 \% no primeiro quarto de notas a 20,1 \% no último quarto.
Entre os ingressantes, esta proporção apresentou o mesmo comportamento : indo de 52,6 \% a 43,3 \%.
A proporção que declarou a parte Fácil/Muito fácil foi maior entre os concluintes ( 13,8 \% ) do que entre os ingressantes ( 6,9 \% ).
Entre os concluintes, 59,4 \% classificaram esta parte com um nível Médio de dificuldade.
A proporção entre os ingressantes foi menor, 44,2 \%.
Entre os concluintes, esta proporção basicamente aumentou com o desempenho, indo de 46,9 \% no primeiro quarto de notas a 66,1 \% no último quarto.
Entre os ingressantes esta proporção também foi crescente, indo de 38,4 \% a 53,7 \%.
Os Gráficos 4.4a e 4.4b apresentam a distribuição das respostas à questão “ Qual o grau de dificuldade desta prova na parte do Componente Específico? ”, segundo a Região onde o curso \textbf{está} situado para os estudantes concluintes e ingressantes que participaram do ENADE/2009 na área de TEATRO.
Cumpre \textbf{lembrar} aqui também que a Região Norte não apresentou nenhum estudante concluinte.
Nenhum ingressante da Região Norte declarou que esta parte foi Muito fácil.
Entre os concluintes, a maior proporção de estudantes que considerou a parte do Componente Específico da prova como de dificuldade média entre todas as Regiões brasileiras, foi na Região Sul com 64,4 \%, já entre os ingressantes, foram os estudantes da Região Norte com 51,4 \%.
A Região com as menores proporções desta resposta foram respectivamente para os concluintes e ingressantes, a Nordeste ( 46,3 \% ) e a Centro-Oeste ( 35,1 \% ). 4.2 Extensão da prova em relação ao tempo total Os Gráficos 4.5a, 4.5b e 4.5c apresentam a distribuição das respostas à questão sobre a extensão da prova em relação ao tempo total, segundo quartos de nota final e grupo de estudantes concluintes e ingressantes que participaram do ENADE/2009 na área de TEATRO.
Tanto para concluintes quanto para ingressantes, a maioria considerou a extensão da prova como adequada em relação ao tempo destinado à sua resolução, respectivamente 62,8 \% e 63,9 \%.
Entre os concluintes, esta proporção cresceu com o desempenho.
A proporção dos que consideraram a extensão da prova muito curta foi maior entre os concluintes ( 1,3 \% ) do que entre os ingressantes ( 0,3 \% ).
A proporção dos que consideraram a extensão da prova curta/muito curta vis-à-vis ao tempo total foi maior entre os concluintes ( 4,4 \% ) do que entre os ingressantes ( 3,9 \% ).
Entre os ingressantes, somente os estudantes com desempenho no terceiro grupo não declararam que a extensão da prova teria sido muito curta com relação ao tempo total.
A maior proporção de concluintes declarando que a prova foi curta/muito curta foi entre os estudantes do primeiro quarto de desempenho, 5,2 \%.
A proporção dos que consideraram a extensão da prova muito longa foi praticamente idêntica entre os concluintes ( 10,8 \% ) e ingressantes ( 10,9 \% ).
A proporção dos que consideraram a extensão da prova longa foi ligeiramente maior entre os concluintes ( 21,9 \% ) do que entre os ingressantes ( 21,3 \% ).
Os Gráficos 4.6a e 4.6b apresentam a distribuição das respostas à questão sobre a extensão da prova em relação ao tempo total, segundo a Região onde o curso \textbf{está} situado para os estudantes concluintes e ingressantes que participaram do ENADE/2009 na área de TEATRO.
Em todas as Regiões brasileiras, a maioria dos estudantes, tanto concluintes como ingressantes, perceberam a prova como adequada.
Os Gráficos 4.6a e 4.6b apresentam a distribuição das respostas à questão sobre a extensão da prova em relação ao tempo total, segundo a Região onde o curso \textbf{está} situado para os estudantes concluintes e ingressantes que participaram do ENADE/2009 na área de TEATRO.
Em todas as Regiões brasileiras, a maioria dos estudantes, tanto concluintes como ingressantes, perceberam a prova como adequada. 4.3 Compreensão dos Enunciados das Questões 4.3.1 Formação Geral Os Gráficos 4.7a, 4.7b e 4.7c apresentam a distribuição das respostas à questão “ Os enunciados das questões da prova na parte de Formação Geral, \textbf{estavam} claros e objetivos? ”, segundo quartos de nota final e grupo de estudantes concluintes e ingressantes que participaram do ENADE/2009 na área de TEATRO.
Com relação aos enunciados das questões da parte de Formação Geral, as opiniões foram positivas : a grande maioria dos concluintes ( 72,7 \% ) e ingressantes ( 74,0 \% ) consideraram claros e objetivos os enunciados de todas ou da maioria das questões.
Entre os concluintes, esta proporção cresceu com o desempenho : 58,3 \% para o primeiro grupo de desempenho e 76,8 \% para o último.
Entre os ingressantes, esta proporção cresceu de igual maneira, indo de 61,1 \% a 86,6 \%.
Em linhas gerais, pode -se dizer que a combinação de “ poucos ” e “ não, nenhum ” decresceu com o desempenho tanto para concluintes como para ingressantes.
Os Gráficos 4.8a e 4.8b apresentam a distribuição das respostas à questão “ Os enunciados das questões da prova na parte de Formação Geral, \textbf{estavam} claros e objetivos? ”, segundo a Região onde o curso \textbf{está} situado para os estudantes concluintes e ingressantes que participaram do ENADE/2009 na área de TEATRO.
Em termos regionais, as proporções de ingressantes e de concluintes que consideraram claros e objetivos os enunciados de todas ou da maioria das questões variaram entre os concluintes de 70,5 \% no Nordeste a 73,6 \% no Sudeste.
Já entre os ingressantes estas proporções variaram de 70,3 \% no Norte a 78,8 \% no Sul.
Com relação às respostas “ poucos ” e “ não, nenhum ” tanto para concluintes como para ingressantes as maiores proporções ocorreram na Região Centro-Oeste, respectivamente 17,5 \% e 16,4 \%. 4.3.2 Componente Específico Os Gráficos 4.9a, 4.9b e 4.9c apresentam a distribuição das respostas à questão “ Os enunciados das questões da prova na parte de formação específica \textbf{estavam} claros e objetivos? ”, segundo quartos de nota final e grupo de estudantes concluintes e ingressantes que participaram do ENADE/2009 na área de TEATRO.
Com relação aos enunciados das questões da parte do Componente Específico, as opiniões foram também positivas : a grande maioria dos concluintes ( 60,0 \% ) e ingressantes ( 63,7 \% ) considerou claros e objetivos os enunciados de todas ou da maioria das questões.
A proporção de concluintes declarando que todos os enunciados \textbf{estavam} claros e objetivos foi decrescente com o desempenho indo de 15,5 \% no primeiro quarto para 11,1 \% para os concluintes com melhor desempenho.
O número de estudantes que declararam que “ não, nenhum ” dos enunciados das questões da parte do Componente Específico era claro e objetivo, foi equivalente com a parte de Formação Geral : 2,9 \% entre os concluintes e 1,0 \% entre os ingressantes.
Os que responderam poucos são 12,9 \% entre os concluintes e 13,6 \% entre os ingressantes.
Os Gráficos 4.10a e 4.10b apresentam a distribuição das respostas à questão “ Os enunciados das questões da prova na parte de formação específica \textbf{estavam} claros e objetivos? ”, segundo a Região onde o curso \textbf{está} situado para os estudantes concluintes e ingressantes que participaram do ENADE/2009 na área de TEATRO.
Em termos regionais, as proporções de estudantes que consideraram claros e objetivos os enunciados de todas ou da maioria das questões variaram entre os concluintes, de 54,2 \% no Centro-Oeste a 65,0 \% no Sul.
Com relação às respostas “ poucos ” e “ não, nenhum ” para concluintes a maior proporção ocorreu na Região Centro-Oeste com 19,8, enquanto para os ingressantes foi a Região Norte com 24,3 \%. 4.4 Suficiência das informações/instruções fornecidas Os Gráficos 4.11a, 4.11b e 4.11c apresentam a distribuição das respostas à questão “ As informações/instruções fornecidas para a resolução das questões foram suficientes para resolvê -las? ”, segundo quartos de nota final e grupo de estudantes concluintes e ingressantes que participaram do ENADE/2009 na área de TEATRO.
Quando \textbf{perguntados} sobre as informações/instruções fornecidas para a resolução das questões, quase 77,0 \% dos concluintes e 77,3 \% dos ingressantes de todo o Brasil responderam que estas eram ou excessivas ou suficientes em todas ou na maioria das questões, demonstrando uma avaliação positiva da prova.
Levando -se em conta o desempenho dos participantes, percebeu -se consideráveis semelhanças entre as opiniões de estudantes, ingressantes e concluintes.
A proporção de estudantes com a opinião supracitada foi crescente com o desempenho entre ingressantes, com proporções que vão de 70,1 \% no primeiro grupo a 89,2 \% no último.
O complemento desta opinião, ou seja, os estudantes que responderam “ somente em algumas ” ou “ em nenhuma delas ”, apresentou uma proporção decrescente com o desempenho para estudantes ingressantes.
A resposta “ não, em nenhuma delas ”, não apareceu apenas entre os ingressantes com o melhor desempenho.
Os Gráficos 4.12a e 4.12b apresentam a distribuição das respostas à questão “ As informações/instruções fornecidas para a resolução das questões foram suficientes para resolvê -las? ”, segundo a Região onde o curso \textbf{está} situado para os estudantes concluintes e ingressantes que participaram do ENADE/2009 na área de TEATRO.
Os Gráficos 4.12a e 4.12b apresentam a distribuição das respostas à questão “ As informações/instruções fornecidas para a resolução das questões foram suficientes para resolvê -las? ”, segundo a Região onde o curso \textbf{está} situado para os estudantes concluintes e ingressantes que participaram do ENADE/2009 na área de TEATRO. 4.5 Dificuldade encontrada ao responder à prova Os Gráficos 4.13a, 4.13b e 4.13c apresentam a distribuição das respostas à questão “ Você se deparou com alguma dificuldade ao responder à prova? ”, segundo quartos de nota final e grupo de estudantes concluintes e ingressantes que participaram do ENADE/2009 na área de TEATRO. \textbf{Perguntados} sobre as dificuldades com as quais se depararam ao responder à prova, 45,3 \% dos ingressantes apontaram como principal dificuldade o desconhecimento do conteúdo, já os 40,3 \% dos concluintes apontaram como principal dificuldade a forma diferente de abordagem.
Na verdade, para a quase totalidade das áreas, ingressantes apresentaram uma proporção de estudantes bem maior do que entre os concluintes declarando como principal dificuldade o desconhecimento da matéria.
O desconhecimento do conteúdo foi apontado somente por 14,4 \% dos concluintes, valores variando conforme o desempenho indo de 22,4 \% no primeiro grupo a 11,1 \% no grupo de melhor desempenho.
A forma diferente de abordagem do conteúdo foi considerada como maior dificuldade por 25,6 \% dos ingressantes, enquanto que a falta de motivação foi a principal dificuldade para 31,4 \% dos concluintes e para 20,7 \% dos ingressantes.
Em termos nacionais, uma menor proporção de ingressantes que de concluintes declarou que não teve qualquer dificuldade para responder à prova, respectivamente 5,0 \% e 8,7 \%.
Os Gráficos 4.14a e 4.14b apresentam a distribuição das respostas à questão “ Você se deparou com alguma dificuldade ao responder à prova? ”, segundo a Região onde o curso \textbf{está} situado para os estudantes concluintes e ingressantes que participaram do ENADE/2009 na área de TEATRO.
Em relação às Regiões brasileiras, a Região Norte não apresentou nenhum estudante concluinte, mas todas as Regiões apresentaram estudantes ingressantes.
O percentual de ingressantes que apontaram o desconhecimento do conteúdo como principal dificuldade ao responder à prova variou de 25,7 \% na Região Norte, a 50,0 \% na Centro-Oeste.
Entre os concluintes, tal percentual variou de 12,0 \% na Região Sudeste a 18,7 \%, na Nordeste.
A forma diferente de abordagem de conteúdo foi citada como maior dificuldade por 25,2 \% dos concluintes na Região Nordeste chegando a 45,0 \% na Região Sudeste.
Já entre os ingressantes os valores foram de 23,4 \% na Região Nordeste a 45,7 \% na Região Norte.
O percentual de concluintes que citaram a falta de motivação como principal dificuldade variou de 25,9 \% na Região Sul a 41,1 \% na Região Nordeste.
Entre os ingressantes as estatísticas correspondentes foram de 16,9 \% na Região Centro-Oeste a 27,9 \%, na Região Nordeste.
Os que declararam que não tiveram qualquer tipo de dificuldade variaram, entre os concluintes, de 6,3 \% na Região Centro-Oeste a 13,0 \% na Região Sul.
Entre os ingressantes este percentual variou de 4,3 \% na Região Sul a 5,7 \% na Região Norte. 4.6 Percepção sobre os conteúdos das questões objetivas da prova Os Gráficos 4.15a, 4.15b e 4615c apresentam a distribuição das respostas à questão “ Considerando apenas as questões objetivas da prova, você percebeu que :... ”, segundo quartos de nota final e grupo de estudantes concluintes e ingressantes que participaram do ENADE/2009 na área de TEATRO.
As respostas possíveis para esta questão \textbf{estavam} relacionadas à \textbf{cobertura} do conteúdo das questões objetivas durante o curso.
A proporção de ingressantes ( 61,2 \% ) que declarou que não estudou ainda a maioria desses conteúdos foi bem maior do que entre os concluintes ( 8,6 \% ).
Esta proporção foi decrescente para os concluintes e ingressantes como função do desempenho.
Por outro lado, a proporção de estudantes que declarou que estudou e aprendeu todos ou a maioria desses conteúdos, entre os ingressantes foi de 21,3 \% e entre os concluintes 68,3 \%.
Esta proporção foi crescente como função do desempenho, tanto para os concluintes como para os ingressantes.
A proporção de estudantes que declarou ter estudado a maioria desses conteúdos, mas não os ter aprendido é bem maior entre os concluintes ( 12,5 \% ) do que os ingressantes ( 5,4 \% ).
Os Gráficos 4.16a e 4.16b apresentam a distribuição das respostas à questão “ Considerando apenas as questões objetivas da prova, você percebeu que :... ”, para os estudantes concluintes e ingressantes que participaram do ENADE/2009 na área de TEATRO, segundo a Região onde o curso \textbf{está} situado.
Em relação às Regiões brasileiras, a Região Norte não apresentou nenhum estudante concluinte ; todas as Regiões apresentaram estudantes ingressantes.
Concluintes que afirmaram não ter estudado ainda a maioria dos conteúdos ocorreram em todas as Regiões, os percentuais variaram de 5,6 \% na Região Sul a 16,0 \%, na Nordeste.
Entre os ingressantes, os percentuais variaram de 45,9 \% na Região Norte a 64,9 \%, na Centro-Oeste.
Em todas as Regiões, a maioria dos concluintes afirmou ter estudado e aprendido todos ou muitos dos conteúdos, com percentuais que variaram de 47,2 \% na Região Nordeste a 80,2 \%, na Sul.
Para os ingressantes, tal percentual variou de 17,3 \%, na Região Centro-Oeste, a 29,7 \%, na Norte. 4.7 Tempo gasto para concluir a prova Os Gráficos 4.17a, 4.17b e 4.17c apresentam a distribuição das respostas à questão “ Qual foi o tempo gasto para concluir a prova? ”, segundo quartos de nota final e grupo de estudantes concluintes e ingressantes que participaram do ENADE/2009 na área de TEATRO.
Quando \textbf{perguntados} quanto ao tempo de conclusão da prova, a maioria respondeu entre uma e duas horas, tanto entre os concluintes ( 39,0 \% ) quanto entre os ingressantes ( 41,1 \% ).
A segunda maior incidência de respostas foi entre duas e três horas, igualmente para ambos os grupos : concluintes ( 34,3 \% ) e ingressantes ( 35,2 \% ).
A proporção daqueles que declararam ter usado as quatro horas, mas não terem conseguido terminar foi bem baixa : 0,9 \% entre os ingressantes e 1,8 \% entre os concluintes.
Os Gráficos 4.18a e 4.18b apresentam a distribuição das respostas à questão “ Qual foi o tempo gasto para concluir a prova? ”, segundo a Região onde o curso \textbf{está} situado para os estudantes concluintes e ingressantes que participaram do ENADE/2009 na área de TEATRO.
A proporção de concluintes que gastaram entre uma e duas horas foi maior na Região Centro-Oeste ( 43,6 \% ) e menor na Região Sudeste ( 36,1 \% ).
Para ingressantes, o maior percentual para este tempo ocorreu na Região Sul ( 46,5 \% ) e o menor na Região Norte ( 31,4 \% ).

% matched lemmas: cada, cobertura, esquema, estar, lembrar, localização, ordem, perguntar
\end{textsample}
